In conclusion, the Arduino-based PID temperature control system effectively
demonstrated progressive improvements in stability and response across the
ON--OFF, P, PI, and PID configurations.
The PI controller achieved the most stable performance, maintaining the target
temperature with the smallest standard deviation and least noise reactivity,
validating its theoretical advantages.
However, the experimentally determined specific heat capacities for aluminum
and brass were significantly higher than accepted values, primarily due to
systematic errors.
These could have included code-related issues such as the chosen constants or
improper power scaling of the PWM signals, circuitry losses through the MOSFET
and wiring, limited thermal insulation, and sensor inaccuracies from the
thermistor.
Had there been more time for the experiment it could have proven beneficial to
incorporate direct voltage and current measurements for accurate power
determination, improved sensor calibration, and refined control code to
eliminate offset bias.
Despite these discrepancies, the experiment successfully demonstrated the
principles and effectiveness of PID-based temperature regulation.